\section{Convergence Bounds of Divergence losses}\label{sec:cv_divergence}

Training losses for Generative flows based on divergence and generalized divergences were introduced in \cite{brunswic2024theory} to show that this wide family of losses is generically unstable.
Further exploration was then conducted \cite{hu2024beyond,silva2024divergence} in the acyclic case. 
Instability of these losses implies that analogues of Theorem \ref{theo:RL_CV_bound} for these losses \textbf{cannot exist}; the following theorem makes this precise.

\begin{theorem}\label{theo:no_bound_divergence}
    Let $\mathcal L(F)$ be a divergence-based FM-loss on a directed graph $\mathcal G$, with
    generator continuous at $1$, depending on a target $\target$ and on a training distribution $\nutrain$; the supremum below ranges over the pairs $(\mathcal G,\target)$ with $\#\mathcal G = N$.
\begin{align*}
    &\forall\, N\geq 2,\quad \forall\, \varepsilon > 0, \\
    &\hspace{1cm}
    \inf_{\nutrain}\; \sup_{\#\mathcal G = N}\; \sup_{\substack{
      F \\ \mathcal L(F) \leq \varepsilon
    }} \; \TV\bigl(s_\tau \| \target\bigr) = 1
\end{align*}
\end{theorem}

To go beyond this limitation one must restrict the setting. Proposition 4.2 of
\citet{silvageneralization} bounds the total variation by an idealized trajectory-balance
residual on finite acyclic graphs of bounded horizon, with constants depending on the state
space and the training distribution --- bounded independently of the trained model on each
graph, but not uniformly in its depth; Appendix
\ref{sec:silva_comparison} restates it with explicit constants and compares it with ours. Such
restrictions fit autoregressive LLM finetuning
\cite{bartoldson2025trajectory,yu2024flow,wang2025gflowgr,ho2024proof,takase2024gflownet},
but a growing trend goes beyond autoregressive generation
\cite{zhu2025llada,nie2025large,sahoo2024simple,cetin2025large,von2025generalized}, where it
does not apply.

\citet{morozov2025revisiting} show that the detailed-balance loss is stable in a weaker sense
on finite graphs when the flow is parameterized with a fixed backward policy. Theorem
\ref{theo:db_stable_frozen} generalizes this quantitatively, from the summable $L^2$-mixing
hypothesis of our universality Theorem \ref{theo:universality_L2}. Before developing the
framework we record a counter-example to the naive generalization. Consider the graph

\begin{equation}\label{eq:counterexample_graph}
\xymatrix{
s_0 \ar[r]& 1  \ar[drrr]\ar[r]& 2  \ar[drr] \ar `ur^l[l]`^d[l] [l]  \ar[r] & \cdots \ar[r]& n \ar[d]\ar[r]& \cdots \ar[r] & 2^n \ar `ur^l[ll]`^d[ll] [ll]  \ar[r] \ar[dll] & \cdots  \\ &&&& s_f&&
}
\end{equation}
on which the uniform backward policy sends backward trajectories to $+\infty$, so the bound of
\citet{morozov2025revisiting}, resting on their expected length, says nothing. A policy that
damps the doubling edge repairs their condition and not ours:
$\BackwardPolicy(j\rightarrow 2j)=\frac{1}{2(j+1)}$ is positive on every edge and, when
backward trajectories enter the graph at states of mean $4.5$, gives $\bar\sigma=9$ exactly,
whereas \emph{every} constant doubling probability diverges --- and
its own mixing sum is infinite all the same.\footnote{Appendix \ref{app:doubling} proves both
halves for the family $\varepsilon(j)=c/(j+1)^s$, of which this policy is the member
$c=\tfrac12$, $s=1$: the $9$ is its closed-form backward length
(Proposition \ref{prop:doubling_length}), and $\widehat B=+\infty$ wherever the family is
positive recurrent (Theorem \ref{theo:doubling_unbounded}). The divergence is priced there on
truncations to $2^{12}$ states.}

\subsection{From Existence to Speed: the Backward Policy Sets the Clock}\label{sub:speed}

Theorem \ref{theo:universality_L2} and, on graphs, Theorem \ref{theo:universality_graphs}
settle \emph{existence}: an ergodic backward policy with summable mixing makes the
parameterization universal, and on a finite graph the frozen-backward family contains an
exactly balanced flow, unique up to scale. They say nothing about how long training takes to
find it. That is what the rest of this section answers, and the answer has a single currency.

Let $G$ be a finite path-connected marked graph, let the frozen backward policy be positive on the edges of the loop closure $\widehat G$ (Definition \ref{def:loop_closure}), and for a state $x$ let $\sigma(x)$ be the time the backward trajectory issued from $x$ spends in the graph, that is, its hitting time of $s_0$. Put
\begin{equation}\label{eq:morozov_variables}
    \sigma_*:=\max_x\mathbb E(\sigma\mid X_0=x),\qquad \bar\sigma:=\mathbb E_{x\sim\widehat\BackwardPolicy(s_f\rightarrow\cdot)}\mathbb E(\sigma\mid X_0=x),
\end{equation}
$\bar\sigma$ being the expected backward-trajectory length of \citet{morozov2025revisiting} (the maximum in $\sigma_*$ running over every state of the loop closure, $s_f$ included), and let $N(x)$ be the expected number of visits to $x$ made by one backward trajectory drawn from the target row $\widehat\BackwardPolicy(s_f\rightarrow\cdot)$, with $N_{\min}:=\min_xN(x)$.

\begin{theorem}[training speed under a frozen backward policy]\label{theo:training_speed}
On a finite marked graph with a frozen backward policy positive on the edges of its loop
closure, and for a training measure of uniformly positive density, gradient training of the
state-space balance loss generated by $(\log x)^2$ --- the generator of the practical DB and
TB losses --- converges to a balanced flow from
\emph{every} positive initialization, first at a universal \L ojasiewicz rate, then
exponentially at a rate fixed by the time backward trajectories spend in the graph. Both
phases hold for the gradient flow and for discrete gradient descent, the latter with a step
size chosen from the initialization --- no step size fixed in advance works from every
initialization. No mixing, spectral-gap or
aperiodicity hypothesis is used for this state-space loss; the DB instance inherits a local
rate through the edge lift wherever the lifted mixing sum is finite (Theorems
\ref{theo:db_stable_frozen} and \ref{theo:local_convergence}).
\end{theorem}
Here $\lambda$ is the backward chain's invariant measure, which is its visit ratio,
\begin{equation}\label{eq:occupation}
    \lambda(x)=\frac{N(x)}{2+\bar\sigma},\qquad\text{so}\qquad \lambda_{\min}=\frac{N_{\min}}{2+\bar\sigma},
\end{equation}
and $\varrho_\sigma$ is the second-phase rate, per unit mass at the crossover:
\begin{equation}\label{eq:morozov_rate}
    \varrho_\sigma\;=\;\frac{g''(1)\,w_{\min}\,\lambda_{\min}}{\sigma_*^{\,2}}\;=\;\frac{2\,w_{\min}\,N_{\min}}{\sigma_*^{\,2}\,(2+\bar\sigma)} ,
\end{equation}
Its dependence on the graph and the policy runs through $\sigma_*$, $\bar\sigma$ and
$N_{\min}$ alone. Full statement, with the
\L ojasiewicz envelope and the explicit crossover time, and proof: Theorem
\ref{theo:training_speed_full}, Appendix \ref{app:divergence}.

The last sentence is the point of the reformulation. What governs training is not how long
the backward chain takes to forget its starting point but how long it takes to \emph{leave the
graph}, and \eqref{eq:occupation} converts one into the other. Three consequences: aperiodicity
leaves the hypotheses, which matters because a leveled --- autoregressive --- graph is periodic
and has infinite mixing sum; the hypothesis becomes checkable, since $\bar\sigma$ and the
$N(x)$ are expectations over backward trajectories of the model one is already training; and
the rate degenerates on the graph
\eqref{eq:counterexample_graph} at every doubling probability --- $\sigma_*$ is a maximum over
states --- and not only in the window where Morozov's bound fails (\S\ref{sec:doubling_cost}). The price is that
$\widehat B_\sigma$ is a conservative envelope: \eqref{eq:morozov_rate} is a guarantee, not a
prediction, and Section \ref{sec:experiments} measures how conservative. Remark
\ref{rem:visit_ratio} separates the regime each constant governs.

\subsection{First Variation of Balance Losses}

Everything rests on the first variation of the $g$-divergence balance loss
$\mathcal L_{g,\nu}(\mu) := \int g\left(\frac{d(\mu T)}{d\mu}\right)d\nu$, for an
\emph{arbitrary} Markov kernel $T$. Four measures meet in its gradient: the background
measure $\nuB$, which carries the parameterization; the weighting measure $\nu$, which is
the training measure $\nutrain$ in practice --- we keep the generic letter because the
statements below hold for any $\nu$ with the positivity each one names and a bounded density,
the latter automatic on a finite state space and not on a continuous one; an invariant measure
$\lambda$ of $T$, the one of the four that the practitioner neither chooses nor samples; and
$\nuG$, any finite reference measure equivalent to $\lambda$, which fixes the geometry the
gradient is read in rather than the objective.

\begin{theorem}\label{theo:first_variation}
For a continuously differentiable generator, the $g$-divergence balance loss of an arbitrary
Markov kernel is differentiable at every flow whose balance ratio is bounded away from $0$ and
$\infty$, and --- as soon as the kernel admits a non-zero finite invariant measure equivalent
to that flow --- its gradient is the reversal of the kernel, minus the ratio, applied to $g'$
of that ratio.
\end{theorem}
The reversal here is conjugation in $L^2(\lambda)$ against the kernel's invariant measure: an
involution sending $T$ to the kernel $T^\lambda$ with $\lambda\otimes T^\lambda = T\otimes\lambda$,
and exchanging a backward policy with the forward policy of the \emph{balanced} flow (Lemma
\ref{lem:adjoint}) --- the forward/backward duality of generative flows. Explicitly,
$\nabla_{\mu}^{\nuG}\mathcal L_{g,\nu} = \frac{d\nuG}{d\lambda}\cdot(T^\lambda-r)\left[g'(r)\frac{d\nu}{d\mu}\lambda\right]$
with $r=\frac{d(\mu T)}{d\mu}$. Full statement and proof: Theorem
\ref{theo:first_variation_full}, Appendix \ref{app:divergence}; the rates below are read at
the invariant measure of the chain in play.

FM, DB and TB are all specializations of it. The DB and TB ratios live on edges and trajectories rather than
states, but they are the same balance loss for the same frozen backward policy lifted to a
power of the state space, so every convergence statement below transfers to them with explicit
constants (Appendix \ref{app:divergence}).

\subsection{Gradient Descent Diffuses the Flow by the Backward Policy}\label{sub:gd_diffusion}

We now linearize the gradient flow of these losses at the balanced solution and read off the convergence speed of training from the mixing of $\BackwardPolicy$. Throughout, $(\Statespace',\lambda',T)$ denotes a chain together with its invariant measure: either $(\Statespace,\lambda,\BackwardPolicy)$ for the FM loss, or its edge lift $(\Statespace^2,\lambda_2,K_2)$ for the DB loss. Then $P$ is the density action of $T$ on $L^2(\lambda')$, $P^\dagger$ its adjoint --- the density action of the reversal $T^{\lambda'}$, by Lemma \ref{lem:adjoint} --- $A := P - I$, and $M_w$ multiplication by $w=\frac{d\nutrain}{d\lambda'}$.

\begin{theorem}[gradient descent as backward diffusion]\label{theo:gd_diffusion}
Near the balanced flow, and for any generator of class $C^3$ near $1$ with non-degenerate curvature there, the
gradient field of a balance loss linearizes to the backward policy's diffusion applied twice, so that training relaxes the current flow guess by diffusing
it through the frozen backward policy --- and for the FM loss with a reversible policy and
uniform training weight ($w\equiv1$) it is exactly the square
of that policy's heat flow, each of its eigenmodes decaying at $g''(1)(1-\beta)^2$.
\end{theorem}
Full statement and proof: Theorem \ref{theo:gd_diffusion_full}, Appendix
\ref{app:divergence}.

The convergence speed of training is therefore a property of the spectrum of that policy.

\begin{theorem}[convergence speed under frozen backward policy]\label{theo:db_stable_frozen}
For a training measure of uniformly positive density, a backward policy with summable
mixing makes gradient descent on the \emph{linearized} FM and DB losses with frozen backward
policy converge to the balanced ray at an explicit exponential rate, provided the step size
stays below an explicit cap; the linearization conserves the total flow normalization along
the way. At vanishing training weight the rate is vacuous.
\end{theorem}
The rate is $\varrho = g''(1)w_{\min}/\widehat B^{\,2}$ with $\widehat B$ the mixing sum,
which for the FM loss is $B$ itself and for DB is $1+B$. Full statement and proof:
Theorem \ref{theo:db_stable_frozen_full}, Appendix \ref{app:divergence}.

\begin{theorem}[local exponential convergence]\label{theo:local_convergence}
On a finite state space, and from an explicit neighbourhood of the unit-mass balanced flow,
the rate above upgrades to the full nonlinear dynamics: halved for the gradient flow, and
quartered for gradient descent below an explicit step.
\end{theorem}
Full statement and proof: Theorem \ref{theo:local_convergence_full}, Appendix
\ref{app:divergence}.

These are statements about the \emph{dynamics}: the balanced ray attracts, exponentially and
at an explicit rate. They are not statements about the loss landscape in the sense of
\citet{brunswic2024theory}, whose Definition 3 calls a loss \emph{stable} when adding a
$0$-flow never decreases it,

\begin{equation}\label{eq:def_stability}
    \mathcal L(F+F_0)\;\geq\;\mathcal L(F) \qquad \text{for every $0$-flow $F_0$ that is a subflow of } F,
\end{equation}
and whose Theorem 3 declares the FM, DB and TB losses unstable. One might hope that freezing
the backward policy repairs \eqref{eq:def_stability}, since it collapses the balanced set to a
single ray. It does not, and the obstruction is the very mechanism that makes training work.

\begin{proposition}[freezing does not restore stability]\label{prop:frozen_unstable}
On every finite path-connected graph carrying a directed cycle that avoids both the source
and the sink, for every admissible generator ($g$ continuous, vanishing at $1$ and positive
elsewhere) and every training measure of positive density, the FM and DB losses with an
everywhere-positive frozen backward policy remain unstable in the sense of
\eqref{eq:def_stability}:
some flow is strictly improved by \emph{adding} a $0$-flow.
\end{proposition}
Full statement and proof: Proposition \ref{prop:frozen_unstable_full}, Appendix
\ref{app:divergence}.

What freezing buys is not
\eqref{eq:def_stability} but the weaker property invoked by \citet{morozov2025revisiting}: the
loss vanishes only on the balanced ray, is coercive transverse to it, and gradient descent
converges there. The three notions --- stability, weak stability, convergence --- are distinct,
and only the last two survive a frozen backward policy.

On trajectories of $\ell$ transitions the TB Hessian at balance is at most that of DB times
$\ell$ and a mixing factor, so a DB step cap divided by both remains valid for TB; no rate is
proved for TB (Proposition \ref{prop:tb_hessian}). Summable mixing, sufficient for the FM
and DB rates, is
the hypothesis of Theorem~\ref{theo:universality_L2} at $p=2$: \emph{the same ergodicity of the backward
policy provides both the universality of the parameterization and the convergence speed of its
training}. Practice, however, often trains the backward policy rather than freezing it
\citep{morozov2025revisiting}; Appendix \ref{app:pbtrain} translates that objective into the
spectral variables used here, where flow regularization becomes a slow drift along the
balanced manifold toward an optimal-transport endpoint, every constant above pays a computable
price along the way, and an entropic term repairs it. Remark \ref{rem:visit_ratio} resolves the speed state by state.

Whether convergence extends beyond the neighbourhood of Theorem \ref{theo:local_convergence} depends on the generator, and the next theorem separates two classes of it: the strictly unimodal generators, and those whose $g'$ vanishes on bands on either side of $1$ --- bands of zero mobility $g''$, on which the diffusion of Theorem \ref{theo:gd_diffusion} freezes. The two classes do not exhaust the generators: bands on one side of $1$ only fall under neither, and within the first, global convergence is proved for $(\log x)^2$ and $(x-1)^2$ only, a bounded strictly unimodal generator not being covered. For $g=(\log x)^2$, the generator of the practical DB and TB losses, one may fear a milder failure on cyclic graphs: a flow stuck in a \emph{distant equilibrium} between the frozen-policy constraint and the circulation-inflating force of Theorem \ref{theo:no_bound_divergence}. The following dichotomy rules this out.

\begin{theorem}[global dichotomy for frozen-backward training]\label{theo:global_dichotomy}
If $g$ is strictly unimodal --- as $(\log x)^2$ and $(x-1)^2$ are --- then every critical point
of the loss is balanced, on every graph and for every training measure with positive density,
and --- for $(\log x)^2$ and $(x-1)^2$ --- on a finite state space the gradient flow
converges globally from every positive initialization. If instead
$g'$ vanishes on bands on either side of $1$, then every flow whose ratios lie in those bands
is frozen, for every training measure; and some such generators freeze an open set of flows
at a definite distance from balance, already on the two-state chain with maximal mixing, so
that global convergence fails for them.
\end{theorem}
Full statement and proof: Theorem \ref{theo:global_dichotomy_full}, Appendix
\ref{app:divergence}.

Convex or real-analytic generators exclude this freezing altogether (Remark
\ref{rem:freezing}).
